% Part: proof-theory
% Chapter: normalization
% Section: introduction

\documentclass[../../../include/open-logic-section]{subfiles}

\begin{document}

\olfileid{pt}{nor}{int}

\olsection{پېژندنه}

که مو يو شرطي فارمول~$!A \lif !B$ !!{prove} کړے وي،
نو ښايي د~\Elim{\lif} په کارولو ئې د~$!B$ په
!!a{proof} کښې وکاروئ۔ د دې دپاره البته د~$!A$
!!a{proof}~$\delta_1$ هم پکار دے۔ د~$!A \lif !B$ ستاسو
!!{proof}~$\delta_3$ ډېر امکان لري چې په~\Intro{\lif}
ختمېږي۔ دا قاعده له پرانيستي فرض~$!A$ څخه د~$!B$
!!a{proof}~$\delta_2$ د~$!A \lif !B$ په ثبوت بدلوي۔ که
همداسې وي، ستاسو وروستے !!{proof} دا بڼه لري:
\[
\AxiomC{$\Discharge{!A}{x}$}
\RightLabel{$\delta_2$}
\DeduceC{$!B$}
\DischargeRule{\Intro{\lif}}{x}
\UnaryInfC{$!A \lif !B$}
\AxiomC{}
\RightLabel{$\delta_1$}
\DeduceC{$!A$}
\RightLabel{\Elim{\lif}}
\BinaryInfC{$!B$}
\DisplayProof
\]
ستاسو تګلاره په دې اغېزمنه ده چې له مخکې موجود يو څيز،
د~$!A \lif !B$ !!a{proof}، مو بيا وکاروۀ، خو خپله
!!{proof} مو اغېزمن نۀ دے۔ د $!A \lif !B$ داخلول او
سمدستي ئې ايستل يو تاو راتاو ښکاري۔ د~$!B$ د !!{prove}
کولو نېغه لار بايد موجوده وي۔
\textbf{سرچينه‌يي سمون OLCMP-167:} د شرطي فارمول ثبوت
د مقدمې له ثبوت څخه بېل نوم لري؛ د سرچينې نقل شوے نوم
د پاييزې پايلې له فرعي ثبوت سره ګډېدۀ۔
\textbf{سرچينه‌يي سمون OLCMP-168:} د رسم ښے فرعي ثبوت
مقدمه ثابتوي، نو د هماغه لومړني ثبوت نښه پکار ده؛ د
شرطي بدنې نقل شوې نښه سمه شوه۔

په حقيقت کښې، نېغه لار تل شته۔ پورته په څېر
``تاو راتاو''، يعنې !!a{introduction} قاعده چې ورپسې
سمدستي !!a{elimination} قاعده وي، د فطري استنتاج له
!!{proof}s څخه تل لرې کېدے شي۔ دې ته \emph{عادي کول}
وايو، او پايله يو \emph{عادي} !!{proof}، يا !!a{proof}
\emph{په عادي بڼه کښې} وي۔

د دې پايلې ثبوت، لکه په سېکوېنټي حساب کښې د پرېکون
لرې کولو قضیه، يو څه پېچلے دے او ډېر حالتونه څېړل غواړي۔
په کلاسيکي \Log{N1c} کښې ئې ثابتول د حدسي \Log{N1i}
په پرتله سخت دي۔ لکه چې د پرېکون لرې کول ښيي چې په
سېکوېنټي حساب کښې د \Cut{} قاعدې په کارولو نور څه
!!{prove} کولے نۀ شو، عادي کول ښيي چې د تاو راتاو په
کارولو د هغه له مخنيوي څخه زيات څه نۀ شو ثابتولے۔ نور
ورته‌والي هم شته: يو عادي !!{proof} د هماغې پايلې له
تر ټولو لنډ تاو راتاو لرونکي !!{proof} څخه ډېر اوږد
کېدے شي؛ لکه چې سېکوېنټي !!{proof}s د پرېکونونو سره د
تر ټولو لنډ بې‌پرېکونه !!{proof} څخه ډېر لنډ کېدے شي۔
د فرعي !!{formula} خاصيت، چې يوه پايله تل په داسې
!!{proof} کښې !!{prove} کولے شو چې يوازې د پايلې او
پاتې پرانيستو فرضونو فرعي !!{formula}s ولري، په حدسي
د فطري استنتاج په نظام کښې له عادي کولو راځي، لکه په
سېکوېنټي حساب کښې چې له پرېکون لرې کولو راځي۔ په لومړي
ترتيب کښې د کميت‌لرونکو فرعي فارمولونو اړوند تعويضي
موردونه هم د همدې خاصيت په معنا کښې شامل دي۔
\textbf{سرچينه‌يي سمون OLCMP-169:} د سرچينې بيان پاتې
پرانيستي فرضونه پرېښي وو۔ مثلاً، د يوه شرطي فرض او د
هغۀ د مقدمې څخه په نېغ ايستلو پايله ثابتېږي، خو شرطي
فرض يوازې د پايلې فرعي فارمول نۀ دے۔ د بيان دقيق حد
پايله او پاتې فرضونه رانغاړي؛ د دې باب رسمي عادي کولو
قضیه حدسي نظام ته ورکړل شوې ده، د کلاسيکي نفي لرونکې
قوي فرعي‌فارمول دعوه دلته نۀ زياتېږي۔

\end{document}
